% ch3 C 金属中的准粒子：费米液体 (原书页 126–131 / PDF p141–p146)

\section{金属中的准粒子：费米液体}

Landau 的“费米液体”（Fermi liquid）理论（73）或许可以被看作“钳位”（clamping）或有效哈密顿量观念的一种问题多得多的应用。Landau 理论试图把准粒子观念应用于相互作用很强的简并费米气体（degenerate Fermi gas），例如金属中的电子，或低温下的液体 He$_3$。

在这里，散射是电子自身之间相互作用的效应，因此若要完全以钳位理论来论证准粒子的使用，就意味着不仅要对离子实的运动加钳位，还要对电子气体本身加钳位。但这件事无法以任何十分严格的方式做到。

Landau 的做法是这样来施加他的“钳位”的：规定准粒子的分布 $n(k,r)$ 应当如何作为金属中位置的函数。然后，正如我们对声子存在下的准粒子所做的那样，他假定一个给定的波包 $q_k^{\dagger}(r)$，其能量随同一区域内其他准粒子的密度而变化：

\begin{equation*}
E_k(r) \;=\; E_k^{0} \;+\; \sum_{k'} f_{kk'}\, n(k', r)
\end{equation*}

有必要保留指标 $k$，用以指明准粒子波包以哪个 $k$ 值为中心，因为低能的准粒子在动量空间中可以铺展到整个费米面。因此我们必须谨慎地处理分布函数 $n$，因为 $k$ 和 $r$ 不是对易的变量，只能粗略地加以指定。

我们至今还没有说明准粒子波包究竟是什么的波包。在这里，即便是系统最低的精确激发态也极其复杂。这一点可以直接从最低阶微扰论看出来。考虑一个粒子 $k$，它已被激发到费米面之上到达 $k_F+a$，这构成无相互作用 Fermi 系统最简单的激发态之一。这一电子可能被散射 $k = k_F + a \to k - q$，此过程同时把第二个电子 $k'$ 激发到 $k'+q$；只要末态中所有的能量都高于费米能，该过程就不为 Pauli 原理所禁止
%FIG{38}: 图 38
（图 38）。毫无疑问各位早已熟知，由此得到的平均自由时间按能量增量的平方变化，

\begin{equation*}
\Delta E \;=\; \frac{\hbar^{2}}{2m}\left[(k_F+a)^{2} - k_F^{2}\right] \;\simeq\; \frac{\hbar^{2}}{m}\, k_F a
\end{equation*}

\begin{equation*}
\frac{\hbar}{\tau} \;=\; \text{consts.} \times \frac{|V|^{2}\,(\Delta E)^{2}}{E_F^{3}}
\end{equation*}

这一点很容易看出，只要注意到：对任何一个空穴留在 $k'$ 态的末态，能量守恒要求

\begin{equation*}
k^{2} - (k-q)^{2} \;=\; (k'+q)^{2} - k'^{2}
\end{equation*}

或

\begin{equation*}
(k - k') \cdot q \;=\; q^{2}
\end{equation*}

这是以 $(k-k')$ 为直径的球面的方程：于是我们知道，$k'+q$ 必定落在以从 $k$ 到 $k'$ 的矢量为直径的球面上。当 $a$ 很小时，这个球面可以用通过 $k$ 和 $k'$ 的平面来近似，因为 $k'$、$k$、$k-q$ 和 $k'+q$ 都必须非常靠近费米面（除了 $k' \simeq -k$ 的情形，而这只占可用态中极小的一部分）。现在很清楚，若 $k'$ 位于费米面以内 $(a\cos\theta)$ 的深度处，则可能存在两类散射过程：若我们选取一个垂直于 $(k-k')$ 且较小的 $q$（$q > |k_F - k'|\cos\theta$，$q < a\cos\theta$），便可以有 $k \to k-q$、$k' \to k'+q$，或者 $k \to k'+q$、$k' \to k-q$（图 39）。无论哪种情形，对给定的 $\theta$，$k'$ 与 $q$ 的可能范围都与 $a$ 成正比，而 $a \propto \Delta E$；于是在这个 $\theta$ 处的总散射概率 $\propto (\Delta E)^{2}$。对 $\theta$ 作积分、从而得到具有给定 $\Delta E$ 的电子的 $\hbar/\tau$ 的完整表达式，这个问题我留作练习。

%FIG{39}: 图 39

使人有理由相信，即便存在散射准粒子概念仍有某种意义的，是这样一个事实：至少在最低阶上，电子的寿命在其能量趋近费米面时非常迅速地变长。如果我们考察更高阶的散射，就会发现不相容原理把它们削减得比我刚才讨论的最低阶效应还要彻底得多。Goldstone、Luttinger、Ward 以及其他若干作者（74）从普遍的场论式、完全重整化（full-renormalization）的进路处理过这个问题。也就是说，他们证明了：通过对微扰论所有阶求和，人们可以如何从无相互作用气体的谱、态密度、以及——把二者一并概括起来的——Green 函数，绝热地延拓到有相互作用气体的相应量。

在绝缘体的情形，各位当还记得，我们把系统的 Green 函数写成

\begin{equation*}
G(k, t) \;=\; \langle g|\, c_k^{\dagger}(t)\, c_k(0)\, |g\rangle
\;=\; Z\, e^{iE_k t} \;+\; \int \mathrm{d}E_n\, f(E_n)\, e^{iE_n t}
\end{equation*}

系统的能谱也可以用以频率表出的 Green 函数来定义：

\begin{align*}
G(k, \omega) \;&=\; i \int G(k, t)\, e^{-i\omega t}\, \mathrm{d}t \\
&=\; \frac{Z}{\omega - E_k} \;+\; \int \frac{f(E)}{\omega - E}\, \mathrm{d}E
\end{align*}

通过 $c_k^{\dagger}$ 产生的那些态的实际能谱，与 $\mathrm{Im}\, G(k)$ 相联系；准粒子态的存在，由 $G$ 在 $E_k$ 处出现一个极点来标志。多体理论所做的，就是证明：至少到微扰论的所有阶，\emph{在金属中}假定 $G$ 在一个复能量

\begin{equation*}
E_k \;+\; i\,\frac{\hbar}{\tau}
\end{equation*}

处有极点是自洽的，于是 $i\hbar/\tau$ 就是能量的虚部，也就是相应谱的宽度。事实证明，这一宽度按 $(\Delta E_k)^{2}$ 趋于零，恰好足以允许一个锐利费米面的存在，因为当准粒子态趋近费米面时，它们的定义便随之变得精确。在金属中，作为时间函数的 Green 函数起初迅速下降，一如它在绝缘体中的表现，然后以频率 $E_k/\hbar$ 振荡；但振荡不再无限延续，而是以 $1/\tau$ 的速率被阻尼（图 40）。

%FIG{40}: 图 40

%FIG{41}: 图 41

相应的谱密度如图 41 所示。同样，正如在绝缘体中那样，我们可以对态 $c_k^{\dagger}\Psi_{g}$ 施加一个能量滤波器，从而由它重构出准粒子态：

\begin{equation*}
q_k^{\dagger}\,\Psi_{g} \;\propto\; \int_{E_k-\frac{\epsilon}{2}}^{E_k+\frac{\epsilon}{2}} c_k^{\dagger}\,\Psi_{g}(E)\,\mathrm{d}E
\qquad \epsilon \;\geq\; \frac{\hbar}{\tau}
\end{equation*}

当然，形式上这很容易写成一个时间积分：

\begin{equation*}
q_k^{\dagger}\,\Psi_{g} \;=\; \int_{0}^{\infty} \mathrm{d}t\;\, e^{\left(i\frac{E_k}{\hbar} \,-\, \frac{\epsilon}{2\hbar}\right)t}\; c_k^{\dagger}(t)\,\Psi_{g}
\end{equation*}

当 $E_k \to 0$ 时，量 $\epsilon$ 可以取得比 $E_k$ 更快地趋于零；这正是我们能够定义一个锐利的费米面、并给某一给定的准粒子态指派一个确定占据概率的原因。

准粒子图像的限度显然是 $\hbar/\tau < E_k \sim k_B T$，而且 \emph{a fortiori}（更不必说）$\hbar/\tau < E_F$；$\hbar/E_F$ 是 $G(t)$ 初始下降的时间常数。

从这一类理论出发，并借助用图形方法证明的各种恒等式，Luttinger 和 Nozières 最近直接从微扰论得到了 Landau 费米液体理论的许多结果（75）。显然，目前还不可能把这种“完全重整化”型的理论重新表述成“有效哈密顿量”或钳位理论。由于“有效哈密顿量”做法更为方便，尤其是对输运问题，这是相对 Landau 直观进路的一个不利之处。此外，微扰论的做法必须假定收敛性，而“微扰论（P.T.）从不收敛”几乎是一句格言；并且可以证明，这一个理论恰恰不收敛。

一个值得提出的好问题是：某种更直截了当的“钳位”式进路是否就不能用于这个问题。我们可以问：正如声子理论中的情形那样，我们能否在不严重改变准粒子性质的前提下，把 $N$ 粒子系统那些必要的自由度钳制住。做到这一点的一种办法，也许是对 $N$ 粒子系统强加一个要求：不允许激发长波长的自由度——也就是说，要求密度涨落

\begin{equation*}
\rho^{Q} \;=\; \sum_{k} c_{k+q}^{\dagger}\, c_{k}
\end{equation*}

或者甚至个别的这类密度涨落

\begin{equation*}
\sum_{\substack{\text{some small}\\ \text{region of } k}} c_{k+q}^{\dagger}\, c_{k}
\end{equation*}

对某个足够小的 $q$ 恒等于零。由于在 $N+1$ 粒子问题中我们可以假定，就长程效应而言准粒子密度与真实粒子密度是相同的（当然仍需附带关于 $q=0$ 极限的那些非常特殊的评注），这将等价于把 Landau 的 $n(k,r)$ 固定下来。

遗憾的是，显然即便如此，也仍不能使准粒子成为系统的精确激发，原因是存在我们曾提到过的“交换”散射过程 $k \to k'+q$、$k' \to k-q$——每一个长波过程都各配有一个这样的过程——此外，还有 $k' \simeq -k$ 那个有趣的区域，在那里所有的 $q$ 都是允许的。也许能找到绕开交换散射问题的某种办法；但耐人寻味的是，微扰论已知的种种失效恰恰与 $k' \simeq -k$ 区域相联系（76）。

令人失望的是，准粒子观念以及费米气体的有效相互作用的完整论证至今尚未出现。由于这些概念在 Landau 准经典分布函数并不是令人满意的进路的那些场合被大量使用——尤其是超导电性和强杂质散射的情形——拥有一个更充分的论证将是极为可取的。

在结束时，我想说：准粒子观念对固体物理是一个无比重要的观念，当然对高能物理也是如此。关于准粒子，各种严格的或或多或少严格的思考方式，给了我们十分宝贵的信心：对物理

% ——上句在原书 p146 末中断于 “...the bases of our physical”，p147 起由后续代理续译衔接。
